%==============================================================================
% ZVF Program -- Pillar 2: THEORY
% "From a Mechanistic Diagnostic to Identifiable, Calibrated, and Optimal
%  Control of the Zero-Variance Fraction in Group-Relative RL"
%
% STATUS: CONDITIONAL THEORY NOTE. T1--T3 are proved under their explicit
% assumptions. T3 audits a declared proxy objective; it does not claim that the
% proxy is the uniquely correct learning objective.
%==============================================================================
\documentclass[11pt]{article}
\usepackage{tikz}
\usetikzlibrary{shapes,arrows.meta,positioning,calc}

%--- Generic, self-contained preamble (no external .sty needed) ----------------
\usepackage[utf8]{inputenc}
\usepackage[T1]{fontenc}
\usepackage[margin=1in]{geometry}
\usepackage{amsmath,amssymb,amsthm,mathtools}
\usepackage{bm}
\usepackage{xcolor}
\usepackage{hyperref}
\usepackage{enumitem}
\usepackage{booktabs}
\usepackage[numbers,sort&compress]{natbib}
\hypersetup{colorlinks=true,linkcolor=blue,citecolor=blue,urlcolor=blue}

%--- Theorem environments ------------------------------------------------------
\theoremstyle{plain}
\newtheorem{theorem}{Theorem}
\newtheorem{lemma}{Lemma}
\newtheorem{proposition}{Proposition}
\newtheorem{corollary}{Corollary}
\theoremstyle{definition}
\newtheorem{definition}{Definition}
\newtheorem{assumption}{Assumption}
\theoremstyle{remark}
\newtheorem{remark}{Remark}

%--- Macros consistent with the existing appendix notation ---------------------
\newcommand{\ZVF}{\ensuremath{\mathrm{ZVF}}}
\newcommand{\GU}{\ensuremath{\mathrm{GU}}}
\newcommand{\Var}{\ensuremath{\mathrm{Var}}}
\newcommand{\Varhat}{\ensuremath{\widehat{\mathrm{Var}}}}
\newcommand{\E}{\mathbb{E}}
\newcommand{\Prob}{\mathbb{P}}
\newcommand{\R}{\mathbb{R}}
\newcommand{\1}{\mathbb{1}}
\newcommand{\Bern}{\mathrm{Bernoulli}}
\newcommand{\indep}{\perp\!\!\!\perp}
\DeclareMathOperator*{\argmax}{arg\,max}
\DeclareMathOperator*{\argmin}{arg\,min}

\title{\textbf{Calibrating the Zero-Variance Fraction in Group-Relative RL}\\[4pt]
\large Estimation, Reliability Budgets, and the Limits of Group-Size Control\\[4pt]
\large ZVF Program, Pillar 2 (Theory)}
\author{Arvind C R\\PES University\\\texttt{arvindcr4@gmail.com}\\[2pt]
\small With AI drafting assistance}
\date{\today}

\begin{document}
\maketitle

\begin{abstract}
The Zero-Variance Fraction (\ZVF{}) is the fraction of prompt groups whose
within-group reward variance is zero in group-relative RL with binary rewards.
It is a useful mechanistic diagnostic, but calibration does not automatically
grant control authority. We separate those questions through three conditional
results. \textbf{T1} treats \ZVF{} correctly as a sample mean of Bernoulli group
indicators and gives binomial-proportion confidence intervals; under curriculum
ordering these intervals cover a local-stage estimand unless batch composition is
stratified. \textbf{T2} converts a population \ZVF{} into a reliability budget:
the number of independent groups required to observe at least one informative
group with probability $1-\delta$. Under the stated assumptions, this bounds
rollouts to a potentially nonzero centered reward-contrast contribution, not
rollouts to policy improvement or a total update. \textbf{T3} audits a
proposed signal-per-rollout objective. Under that declared objective its optimum
collapses to the prior-independent tie $G\in\{2,3\}$; it therefore supplies no
data-adaptive group-size set-point. A richer variance-, Fisher-, wall-clock-, or
outcome-aware objective is required before theory can justify adaptive $G$.
Checked-in sampling experiments validate the operational accounting, while every
assumption and scope boundary is stated inline. The result is a calibrated sensor and a
conservative reliability trigger---not a proof that a particular controller or
group size improves held-out performance.
\end{abstract}
\paragraph{August 2026 evidence boundary.} Theoretical homogeneous-reward
statements concern centered reward contrast and do not assert zero total
gradient, canonical-GRPO equivalence, or an intervention outcome.

%==============================================================================
\section{Introduction}
%==============================================================================
The ZVF Program's Pillar~1 contributed a \emph{mechanism}: in group-relative,
critic-free RL (GRPO and relatives \cite{shao2024deepseekmath}), the submitted
runner's centered reward-contrast term depends on the within-group spread of rewards.
When all $K$ rollouts of a group collapse to the same binary outcome, the
group-relative advantage is identically zero and that centered term is zero under
the stated assumptions; this does not assert a zero total gradient. The \ZVF{}
counts the fraction of such degenerate groups in a batch.
Pillar~1 made this definition rigorous, showed it is \emph{not} a tautological
re-expression of the mean reward, and validated it as a portable early-warning
signal \cite{zvfaudit2026}.

A mechanistic diagnostic, however, is a measurement, and a measurement is only
as good as its calibration, its predictive content, and its actionability.
Three gaps remain before \ZVF{} can \emph{drive} a training-time controller
rather than merely annotate a training curve:

\begin{enumerate}[leftmargin=1.5em,itemsep=2pt]
  \item \textbf{Identifiability / calibration.} The reported $\ZVF_t$ is an
        estimate from finitely many groups. Two runs reporting
        $\ZVF_t = 0.86$ and $0.81$ may be statistically indistinguishable. We
        need its sampling distribution and a confidence interval, or any
        threshold $\tau_{\ZVF}$ is a point estimate masquerading as a decision.
  \item \textbf{A reliability cost.} The operational claim of Pillar~1 is local
        and mechanical: at $\ZVF_t = 1$ the centered reward-contrast term is
        zero under the stated GRPO assumptions; this does not rule out other
        total-gradient contributions. We strengthen this
        into a \emph{finite-sample} statement: high \ZVF{} raises the budget
        required to observe at least one informative group with stated reliability.
  \item \textbf{Limits of action.} If high \ZVF{} wastes rollouts, can a
        signal-per-rollout objective identify the right group size? We audit one
        natural proxy and show that it does \emph{not} yield a data-adaptive
        set-point.
\end{enumerate}

This paper supplies one theorem per gap (T1--T3). We deliberately keep
notation aligned with the Pillar~1 appendix and with the U-statistic analysis
of \citet{zhou2026demystifyinggrpo}. The mixed-group probability used below is
the elementary Bernoulli identity $1-p^G-(1-p)^G$, so it requires no separate
attribution. This alignment lets the three results compose with existing work
rather than re-deriving it.

\paragraph{Related work: interventions without a measurement theory.}
The phenomenon \ZVF{} measures is no longer a niche observation. The TMLR
survey of agentic RL \cite{agenticrlsurvey2026} catalogs, in a single table,
more than twenty GRPO-family variants, and a striking fraction of them
intervene directly on dead groups: DAPO's dynamic sampling drops
zero-gradient groups online \cite{yu2025dapo}; GRESO filters prompts
\emph{before} rollout when they are predicted to produce single-outcome
groups \cite{greso2025}; EDGE-GRPO reshapes advantages explicitly to
mitigate what its authors --- and now the survey --- call \emph{advantage
collapse} \cite{edgegrpo2025}; DARS reallocates rollout budget from
medium-difficulty prompts to the hardest ones \cite{dars2025}. We adopt
\emph{advantage collapse} as the field's name for the phenomenon and read
this convergence as evidence for the premise of the ZVF Program: the
dead-group fraction is a first-order quantity in critic-free RL. Our
position, stated plainly, is that these methods intervene on related but not
identical signals; not all estimate batch-level ZVF. T1 supplies a sampling
model for the explicitly defined ZVF estimator, and T2 gives a conditional
reliability budget. T3 audits one proxy and proves its fixed optimum
$G\in\{2,3\}$; it supplies no adaptive set-point. The results can inform a
measurement layer, but do not prove that every cited method needs this
particular estimator or objective.

\paragraph{Honesty contract.} T1--T3 are conditional mathematical statements:
their proofs are complete under the assumptions stated below. This does not make
the assumptions empirically universal, and T3 does not elevate its declared
proxy into a uniquely correct learning objective. The companion
\texttt{THEORY\_NOTES.md} records the proof audit and the empirical extensions
that remain outside the theorems.

%==============================================================================
\section{Preliminaries and Notation}\label{sec:prelim}
%==============================================================================
We adopt the Pillar~1 setup verbatim where possible.

\paragraph{Batch and groups.} At optimizer step $t$, a GRPO batch $B_t$
consists of $|B_t| = m$ prompt-groups. Each group $g$ corresponds to a single
prompt $x_g$ and contains $K$ rollouts (the \emph{group size}; we write $G$
when treating it as a control variable in Section~\ref{sec:T3}) with scalar
outcome rewards $r_{g,1},\dots,r_{g,K} \in \R$.

\paragraph{Binary verifiable rewards.} Throughout T1--T3 we take rewards
binary, $r_{g,k}\in\{0,1\}$, as produced by a verifier (correct / incorrect).
Conditional on the prompt $x_g$ and the current policy $\pi$, rollouts within a
group are i.i.d.\ Bernoulli with prompt-specific success probability
\[
p_g \;\coloneqq\; \Prob_{\pi}\!\big(r=1 \mid x_g\big) \in [0,1].
\]

\paragraph{Within-group sample variance and the zero-variance event.} Let
$\Varhat_g \coloneqq \frac{1}{K-1}\sum_{k=1}^{K}(r_{g,k}-\bar r_g)^2$ be the
unbiased within-group variance, $\bar r_g = \frac1K\sum_k r_{g,k}$. For binary
rewards, $\Varhat_g = 0$ iff all $K$ rollouts share the same outcome. With a
numerical tolerance $\varepsilon$ (e.g.\ $10^{-6}$ for $[0,1]$ rewards) the
\emph{zero-variance indicator} of group $g$ is
\begin{equation}\label{eq:zv-indicator}
Z_g \;\coloneqq\; \1\!\left[\Varhat_g \le \varepsilon\right]
\;\overset{\text{binary } r}{=}\; \1\!\left[\textstyle\sum_k r_{g,k} \in \{0,K\}\right].
\end{equation}
The equality requires $K\ge2$ and $0\le\varepsilon<1/K$, since the
smallest positive binary sample variance is $1/K$. We use $K$ for a fixed
measurement size and $G$ for the candidate size optimized in T3.

\begin{definition}[Zero-Variance Fraction, Pillar~1]\label{def:zvf}
\[
\ZVF_t \;=\; \frac{1}{m}\sum_{g\in B_t} Z_g .
\]
\end{definition}

\paragraph{Gradient Utilization (GU).} Following the program-wide convention
we set $\GU_t \coloneqq 1-\ZVF_t$, the fraction of groups carrying usable
signal; $\ZVF_t$ itself is the batch-level uniformity rate, and we reserve
the per-group, population-level object for analysis. For binary rewards the population zero-variance probability of a
prompt with success rate $p$ at group size $K$ is the \emph{kernel}
\begin{equation}\label{eq:kernel}
h_K(p) \;\coloneqq\; \Prob\big(Z_g = 1 \mid p_g=p\big) \;=\; p^{K} + (1-p)^{K},
\end{equation}
which is exactly the per-prompt term in the Pillar~1 null-model identity
$\E[\ZVF_t] = \E_{p_g}[\,p_g^{K} + (1-p_g)^{K}\,]$.

\paragraph{Reward-density prior.} For T3 we posit a prior density $\phi$ over
per-prompt success probabilities, $p_g \overset{\text{i.i.d.}}{\sim}\phi$ on
$[0,1]$ (the ``reward-density prior''), e.g.\ a Beta$(\alpha,\beta)$ family.
This is the population-level analogue of the empirical difficulty distribution
$\{p_g\}$ used in Pillar~1.

\paragraph{Standing assumptions.} We collect the assumptions the three results
lean on; each theorem cites the subset it needs.
\begin{assumption}[Binary verifiable rewards]\label{ass:binary}
$r_{g,k}\in\{0,1\}$, produced by a deterministic verifier on policy samples.
\end{assumption}
\begin{assumption}[Within-group i.i.d.]\label{ass:within-iid}
Conditional on $x_g$ and $\pi$, the $K$ rollouts of group $g$ are i.i.d.\
$\Bern(p_g)$.
\end{assumption}
\begin{assumption}[Across-group i.i.d.\ prompts]\label{ass:across-iid}
The prompts (hence $p_1,\dots,p_m$) in a batch are i.i.d.\ draws from a fixed
population distribution. Curriculum sampling, replay, and within-epoch
correlation violate this assumption; in those settings T1 estimates a local
stage quantity unless the design or variance estimator is changed.
\end{assumption}
\begin{assumption}[Non-degenerate support]\label{ass:nondegen}
The reward-density prior $\phi$ assigns positive mass to $(0,1)$, so that
$\E_\phi[h_K(p)] \in (0,1)$ and the population \ZVF{} is interior (rules out
the saturated and trivial regimes flagged in Pillar~1's scope section).
\end{assumption}

%==============================================================================
\section{T1 --- ZVF as a Calibrated Binomial-Proportion Estimator}\label{sec:T1}
%==============================================================================
We treat $\ZVF_t$ as the sample mean of i.i.d.\ bounded group indicators and
derive its asymptotic normality via the binomial-proportion CLT. The result is
a closed-form $95\%$ confidence interval per training step. (We do \emph{not}
claim genuinely U-statistic content for the batch-level \ZVF{}; an order-1
U-statistic is just a sample mean. A degree-2 U-statistic would be needed only
if we estimated a second-order functional such as $\Var(Z_g)$ across groups,
which is not required for the per-step CI.)

\paragraph{The estimand.} Define the population \ZVF{} at group size $K$ as
\begin{equation}\label{eq:theta}
\theta_K \;\coloneqq\; \E_{p\sim\phi}\big[h_K(p)\big]
\;=\; \E_{p\sim\phi}\big[p^{K}+(1-p)^{K}\big] \;=\; \E[Z_g].
\end{equation}
By Definition~\ref{def:zvf}, $\ZVF_t = \frac1m\sum_{g} Z_g$ is the natural
plug-in estimator $\widehat\theta_K$ of $\theta_K$.

\begin{lemma}[Unbiasedness]\label{lem:ustat}
Under Assumptions~\ref{ass:binary}--\ref{ass:across-iid}, $\ZVF_t$ is
unbiased for $\theta_K$, i.e. $\E[\ZVF_t]=\theta_K$ exactly (finite-sample
unbiased) for every $m\ge 1$.
\end{lemma}
\begin{proof}
$\ZVF_t = \frac1m\sum_g Z_g$ is the sample mean of the bounded i.i.d.\
indicators $Z_g\in\{0,1\}$; unbiasedness is linearity of expectation
together with $\E[Z_g]=h_K(p_g)$ averaged over $p_g\sim\phi$, giving
\eqref{eq:theta}. \textbf{We do NOT claim ZVF is a U-statistic of
order $\ge 2$: the U-statistic framing is a category error here (an
order-1 U-statistic with symmetric kernel is a sample mean, so the
asymptotic normality is the trivial binomial-proportion CLT, not
Hoeffding's U-statistic CLT).} The genuine U-statistic content would
require a degree-2 kernel, which is not the estimand in this paper. No
U-statistic claim is needed for the per-step ZVF trajectory.
\end{proof}

\begin{theorem}[Asymptotic normality and closed-form CI for \ZVF{}]
\label{thm:T1}
Fix group size $K$. Under Assumptions~\ref{ass:binary}--\ref{ass:across-iid}
and \ref{ass:nondegen}, as the number of groups $m\to\infty$,
\begin{equation}\label{eq:clt}
\sqrt{m}\,\big(\ZVF_t - \theta_K\big)
\;\xrightarrow{\;d\;}\; \mathcal{N}\!\big(0,\ \sigma_K^2\big),
\qquad
\sigma_K^2 = \theta_K(1-\theta_K),
\end{equation}
because $Z_g\in\{0,1\}$ are i.i.d.\ Bernoulli$(\theta_K)$. Consequently a
closed-form asymptotic $95\%$ confidence interval for the population \ZVF{} at
step $t$ is
\begin{equation}\label{eq:ci}
\boxed{\;
\ZVF_t \;\pm\; 1.96\,\sqrt{\frac{\ZVF_t\,(1-\ZVF_t)}{m}}\;}
\end{equation}
and a Wilson-score interval is recommended near the boundaries
$\ZVF_t\in\{0,1\}$ where the Wald interval \eqref{eq:ci} degenerates.

\paragraph{Empirical calibration (E-T1).} The earlier 350-prompt analysis is
superseded for numerical reporting by the complete 512-prompt pool in
Section~\ref{sec:empirical}. That section reports shuffled Wald coverage
0.93--0.96 at $m\ge64$, Wilson coverage 0.95--0.98, sorted global coverage
0.07--0.22, local coverage 0.944, and stratified coverage 0.996.
These are frozen-pool resampling results, not universal coverage guarantees.
Wilson is the operational interval; the theorem above is asymptotic.
\end{theorem}
\begin{proof}
Each $Z_g$ is a Bernoulli$(\theta_K)$ variable (its mean over the within-group
randomness \emph{and} the prompt draw is $\theta_K$). Under
Assumption~\ref{ass:across-iid} the $\{Z_g\}_{g=1}^m$ are i.i.d., so the
Lindeberg--L\'evy CLT gives \eqref{eq:clt} with variance
$\theta_K(1-\theta_K)$; Slutsky with the consistent plug-in
$\widehat\sigma_K^2=\ZVF_t(1-\ZVF_t)$ yields the Wald CI \eqref{eq:ci}.
The Wilson interval is the standard small-$m$/boundary correction.
Because $Z_g$ is binary, the two-stage randomness (draw $p_g\sim\phi$, then
$K$ Bernoulli trials) collapses to a single Bernoulli marginal:
$\Var(Z_g)=\theta_K(1-\theta_K)$ exactly. The only subtlety is that the
reporting pipeline must average the group indicators $Z_g$, not the conditional
probabilities $h_K(p_g)$; Definition~\ref{def:zvf} does the former.
\end{proof}

\begin{remark}[Connection to ``Demystifying GRPO'' and mixed-group probability]
\label{rem:connections}
The kernel $h_K(p)=p^K+(1-p)^K$ is the complement of the \emph{mixed-group
probability} $1-h_K(p)$, an elementary consequence of independent Bernoulli
trials. T1 turns that per-prompt probability into a calibrated batch-level
estimand with a CI. \citet{zhou2026demystifyinggrpo} instead analyze the GRPO
policy gradient as a U-statistic and derive group-size scaling results. The
present Bernoulli estimand is complementary: it quantifies whether a sampled
group supplies any reward contrast, not the full policy-gradient distribution.
No stronger equivalence between the two analyses is claimed here.
\end{remark}

%==============================================================================
\section{T2 --- A Wasted-Compute Lower Bound}\label{sec:T2}
%==============================================================================
We now show that a high population \ZVF{} yields a reliability budget for
observing a non-degenerate reward-gradient event. The result does not imply a
policy-improvement guarantee.

\paragraph{What ``improvement is possible'' means.} A GRPO update at step $t$
can contribute a nonzero centered reward-gradient term under the stated
assumptions only if at least one group in the batch is
\emph{informative}, i.e.\ has $Z_g=0$ (non-uniform outcomes). Define the
per-prompt informativeness probability $q_g \coloneqq 1-h_K(p_g) =
1-p_g^K-(1-p_g)^K$, the population fraction of informative groups
$\bar q \coloneqq 1-\theta_K = \E_\phi[q_g] = 1-\E[\ZVF_t]$.
The population parameter $\theta_K$ is distinct from its random batch estimate
$\ZVF_t$; substituting the latter does not preserve the guarantees below.

\begin{lemma}[Rollouts per informative group]\label{lem:rollouts-per-group}
Under Assumptions~\ref{ass:binary}--\ref{ass:nondegen}, drawing groups of
size $K\ge2$ until the first informative group ($Z_g=0$) requires, in
expectation, $1/\bar q$ groups, i.e.\ $K/\bar q$ rollouts:
\begin{equation}\label{eq:exp-rollouts}
\E[\text{rollouts to first informative group}] \;=\; \frac{K}{1-\theta_K}.
\end{equation}
\end{lemma}
\begin{proof}
For each new group, Assumption~\ref{ass:across-iid} draws a fresh
$p_g\sim\phi$. Marginalizing that draw gives
$\Prob(Z_g=0)=\E_\phi[1-h_K(p_g)]=\bar q$. Independence of the prompt draws and
rollouts makes the indicators $\{1-Z_g\}$ i.i.d.\ Bernoulli$(\bar q)$, even
though their conditional success probabilities vary with $p_g$. Therefore the
number of groups through the first informative group is geometric with mean
$1/\bar q$; multiplying by $K$ proves \eqref{eq:exp-rollouts}. If instead a
latent difficulty is held fixed across several groups, the i.i.d.-mixture
assumption fails and this equality is not claimed.
\end{proof}

\begin{theorem}[Wasted-compute reliability bound]\label{thm:T2}
Fix group size $K\ge2$ and $\delta\in(0,1)$. Under
Assumptions~\ref{ass:binary}--\ref{ass:nondegen}, let $\theta_K\in(0,1)$
be the population zero-variance probability. Obtaining at least one
informative group with probability at least $1-\delta$ requires and is
ensured by an integer number of independent groups satisfying
\begin{equation}\label{eq:T2-groups}
n_g\ge\left\lceil\frac{\ln(1/\delta)}{-\ln\theta_K}\right\rceil,
\qquad
N_{\mathrm{disc}}(\theta_K)\coloneqq
K\left\lceil\frac{\ln(1/\delta)}{-\ln\theta_K}\right\rceil.
\end{equation}
Here $N_{\mathrm{disc}}$ is the exact rollout quantile. Its continuous lower
bound is
\[
N_{\mathrm{cont}}(\theta_K)\coloneqq
K\frac{\ln(1/\delta)}{-\ln\theta_K}\le N_{\mathrm{disc}}(\theta_K).
\]
A rollout allocation $B$ permits $\lfloor B/K\rfloor$ complete groups and
meets the target iff $B\ge N_{\mathrm{disc}}(\theta_K)$; exceeding only
$N_{\mathrm{cont}}$ is insufficient. Both quantities diverge as
$\theta_K\to1^-$. This is not a bound on policy improvement or total gradients.
\end{theorem}
\begin{proof}
With i.i.d.\ groups, all $n_g$ groups are uninformative with probability
$\theta_K^{n_g}$. Requiring this probability to be at most $\delta$ gives
$n_g\ge\ln(1/\delta)/(-\ln\theta_K)$. Rounding up to a whole group and
multiplying by $K$ proves \eqref{eq:T2-groups}. The divergence follows from
$-\ln\theta_K\to0^+$; more precisely,
$N_{\mathrm{cont}}(\theta_K)\sim K\ln(1/\delta)/(1-\theta_K)$.
An informative group can supply a potentially nonzero centered reward-contrast
contribution, but the proof uses no total-update or policy-improvement claim.
Under correlated groups the product probability $\theta_K^{n_g}$ need not
hold. Corollary~\ref{cor:T2-traj} handles estimation of the population rate.
\end{proof}

\begin{corollary}[Certified compute floor from an observed trajectory]
\label{cor:T2-traj}
Under the assumptions of T2, let $X=m\ZVF_t$ from a diagnostic batch of
$m\ge1$ groups, and fix $\alpha\in(0,1)$. The one-sided
Clopper--Pearson lower confidence bound is
\[
L_\alpha(X;m)=
\begin{cases}
0, & X=0,\\
I^{-1}_{\alpha}(X,m-X+1), & X>0,
\end{cases}
\]
where $I^{-1}$ is the inverse regularized incomplete Beta function. Then
$\Prob_{\theta_K}\{\theta_K\ge L_\alpha(X;m)\}\ge1-\alpha$.
Extend $N_{\mathrm{cont}}(0)=0$ by continuity and
$N_{\mathrm{disc}}(0)=K$ (one group suffices at zero population failure
probability). Both functions are nondecreasing, so for
$j\in\{\mathrm{cont},\mathrm{disc}\}$,
\[
\Prob_{\theta_K}\!\left\{N_j(\theta_K)\ge
N_j\big(L_\alpha(X;m)\big)\right\}\ge1-\alpha.
\]
Thus $N_{\mathrm{disc}}(L_{0.05}(X;m))$ is a conservative 95\% lower
confidence bound on the required rollout budget, not a sufficient allocation.
At $X=0$ it certifies only the trivial one-group minimum; at $X=m$,
$L_\alpha(m;m)=\alpha^{1/m}<1$, so this lower budget bound stays finite.
Wilson bounds remain an operational approximation, not the basis of this
finite-sample certificate.

Conversely, a sufficient budget certificate uses an \emph{upper} bound
$U_\alpha(X;m)=I^{-1}_{1-\alpha}(X+1,m-X)$ for $X<m$, and
$U_\alpha(m;m)=1$. With $N_{\mathrm{disc}}(1)=+\infty$,
\[
\Prob_{\theta_K}\!\left\{N_{\mathrm{disc}}(\theta_K)\le
N_{\mathrm{disc}}\big(U_\alpha(X;m)\big)\right\}\ge1-\alpha.
\]
This is a confidence statement over diagnostic samples about a fresh batch's
$1-\delta$ reliability target, not a claim that confidence $1-\alpha$ and
future-batch success probability $1-\delta$ are the same quantity. These
are pointwise guarantees for a fixed diagnostic sample size; repeated online
checks require separate sequential or multiplicity control.
\end{corollary}

%==============================================================================
\section{T3 --- Auditing a Signal-per-Rollout Group-Size Objective}\label{sec:T3}
%==============================================================================
T2 says high \ZVF{} raises the reliability budget but does not say what action
improves learning. T3 audits a natural signal-per-rollout proxy. Its failure to
produce a data-dependent optimum is a result: this objective cannot supply an
adaptive controller set-point.

\paragraph{Learning signal per group.} For binary rewards at group size $G$,
we declare a diagnostic proxy: weight the mixed-group probability by the
external prompt-difficulty weight $p(1-p)$. To distinguish this choice from
realized sample variance, let $C\sim\mathrm{Binomial}(G,p)$ and
$V=C(G-C)/(G(G-1))$. Homogeneous groups already have $V=0$, hence
\[
\E[V\1\{Z_g=0\}\mid p]=\E[V\mid p]=p(1-p).
\]
The exact expected sample variance per rollout, averaged over $\phi$, is
$\E_\phi[p(1-p)]/G$ and has the unique integer optimum $G=2$ under
Assumption~\ref{ass:nondegen}. The extra mixed-group factor below instead
creates a composite heuristic; its $\{2,3\}$ tie is not a sample-variance
identity or a policy-gradient optimum:
\begin{equation}\label{eq:signal-per-group}
S(p,G)\coloneqq (1-h_G(p))p(1-p),\qquad h_G(p)=p^G+(1-p)^G.
\end{equation}

\paragraph{Signal per rollout.} Each group of size $G$ costs $G$ rollouts, so
the efficiency objective (expected signal per unit sampling cost, averaged over
the prior) is
\begin{equation}\label{eq:efficiency}
J(G) \;\coloneqq\; \frac{1}{G}\,
\E_{p\sim\phi}\!\big[\,(1-h_G(p))\,p(1-p)\,\big].
\end{equation}

\begin{theorem}[Universal discrete optimum of the declared proxy]\label{thm:T3}
Under Assumptions~\ref{ass:binary},~\ref{ass:within-iid},~\ref{ass:nondegen},
the population objective in \eqref{eq:efficiency} satisfies
\begin{equation}\label{eq:J23}
J(2)=J(3)=\E_{p\sim\phi}\big[(p(1-p))^2\big]>J(G)
\quad\text{for every integer }G\ge4.
\end{equation}
Consequently,
\begin{equation}\label{eq:Gstar}
\argmax_{G\in\{2,3,\dots\}} J(G)=\{2,3\},
\end{equation}
independently of $\phi$. This proxy supplies no data-adaptive integer set-point.
\end{theorem}
\begin{proof}
Let $C\sim\mathrm{Binomial}(G,p)$ count successes in a group. The mixed-group
event is $M=\{1\le C\le G-1\}$ and
$\Prob(M)=1-p^G-(1-p)^G$. On $M$,
$C(G-C)\ge G-1$, with equality for every mixed outcome when $G\in\{2,3\}$.
Moreover,
\[
\E[C(G-C)]=G(G-1)p(1-p).
\]
Therefore
\[
\Prob(M)\le\frac{\E[C(G-C)]}{G-1}=G p(1-p).
\]
For $G=2$ and $G=3$ equality holds pointwise. For $G\ge4$ and
$p\in(0,1)$, outcomes with $2\le C\le G-2$ have positive probability and
strictly exceed the lower bound $G-1$, so the inequality is strict. Multiplying
by $p(1-p)/G$ and integrating over $\phi$ proves
$J(2)=J(3)>J(G)$ for every integer $G\ge4$ under
Assumption~\ref{ass:nondegen}.
\end{proof}

\begin{remark}[Continuous relaxation is not an integer control rule]
For real $G>0$, $J$ is differentiable: the terms $p^G|\ln p|$ and
$(1-p)^G|\ln(1-p)|$ are uniformly bounded by $1/(eG)$, justifying
differentiation under $\phi$ locally in $G$. An interior stationary point
therefore satisfies
\begin{equation}\label{eq:foc}
\E_\phi[(1-h_G(p))p(1-p)]
=-G\E_\phi[(\partial_G h_G(p))p(1-p)],
\end{equation}
where $\partial_G h_G(p)=p^G\ln p+(1-p)^G\ln(1-p)$, with endpoint
terms defined by continuity. Since $J(2)=J(3)$, Rolle's theorem ensures
at least one stationary point between them. This does not change the exact
integer maximizer set proved above, even for a point-mass prior.
\end{remark}

\begin{corollary}[Beta-prior closed form]\label{cor:beta}
If $\phi=\mathrm{Beta}(\alpha,\beta)$, then
\begin{align}
J(G)=\frac{1}{G\,B(\alpha,\beta)}\big[&B(\alpha+1,\beta+1)
-B(\alpha+G+1,\beta+1)\nonumber\\
&-B(\alpha+1,\beta+G+1)\big].
\end{align}
Thus $J(G)$ and its continuous derivative are available in closed form.
Given \eqref{eq:J23}, however, the Beta closed form is informative for the
\emph{value} of $J$ but not for the argmax, which is $\{2,3\}$ for every
$(\alpha,\beta)$.
\end{corollary}

%==============================================================================
\section{Connection to the Pillar~1 Adaptive Controller}\label{sec:controller}
%==============================================================================
\noindent\emph{Scope note (2026-07-11 canonicalization).} This paper is
canonical for the theorem-level treatment (T1--T3 and their proof gaps);
the controller's engineering --- trigger design, hysteresis, multi-seed
validation, and design rules --- is treated in the companion Pillar-7
paper \emph{``From Diagnostic to Controller: A Signal-Starvation Theory of
GRPO and an Adaptive Group-Size Intervention''}
(\texttt{platform\_hybrid/paper/paper\_P7\_zvf\_controller}).

The results suggest a restricted prospective interface, not a proved online
control law. Its sampling assumptions must hold for the next batch:
\begin{enumerate}[leftmargin=1.5em,itemsep=2pt]
  \item \textbf{Measure (T1).} From a separate diagnostic batch at fixed $K$,
        compute $\ZVF_t$ and the one-sided Clopper--Pearson lower bound in
        Corollary~\ref{cor:T2-traj}. Wilson is an operational approximation;
        the displayed two-sided Wald interval is not this exact certificate.
  \item \textbf{Decide (T2).} Write $L=L_\alpha(m\ZVF_t;m)$. If
        $N_{\mathrm{disc}}(L)>B$ for the next-batch rollout allocation $B$,
        that fresh batch cannot meet the stated reliability target, with the
        certificate's confidence, under an unchanged population and fixed $K$.
        This certifies \emph{insufficient budget}, not a realized batch failure.
        The converse $N_{\mathrm{disc}}(L)\le B$ certifies no sufficiency;
        that requires the upper-bound construction in
        Corollary~\ref{cor:T2-traj}. No retrospective probability claim is
        made about the already observed diagnostic batch.
  \item \textbf{Act (T3, revised).} Under objective \eqref{eq:efficiency}
        the per-rollout optimum is the universal $G^\star\in\{2,3\}$
        (eq.~\ref{eq:J23}), so the controller cannot derive a
        \emph{data-dependent} group size from T3 as stated; re-estimating
        $\widehat\phi_t$ adds nothing. The escalation policy must instead
        be separately justified: a rate measured at one $K$ does not generally
        identify $\theta_{K'}$ at another size. T2 alone cannot license
        escalation without a cross-size model or new measurements.
\end{enumerate}
At a fixed number of fresh groups $m$, raising $K$ reduces
$\theta_K^m$ but spends more rollouts. At a fixed rollout allocation $B$,
the failure probability is instead $\theta_K^{\lfloor B/K\rfloor}$ and
can increase. For an analytic counterexample (not measured training data),
let half the prompts have $p=0$ and half $p=1/2$. At $B=8$, four groups of
size 2 fail with probability $(3/4)^4=81/256$, whereas one group of size 8
fails with probability $129/256$. Group-size escalation therefore needs an
explicit cost convention and cannot be justified from decreasing $\theta_K$
alone.
This is the formal interface between Pillar~2 (theory) and Pillar~1 (the
mechanistic controller): T1 supplies \emph{sensor calibration}, T2 supplies a
\emph{reliability trigger}, and T3 rules out the proposed proxy as a
data-adaptive actuator set-point. Any escalation rule remains an empirical
policy until a richer objective and a prospective matched-budget test support it.
The obsolete recursion $G^\star\leftarrow\widehat\phi\leftarrow\mathrm{data}$
is not part of this controller because T3 proves that the proxy optimum is
prior-independent. The remaining finite-action escalation rule is an empirical
policy with hysteresis; its closed-loop performance is tested prospectively in
Pillar~7 rather than asserted as a theorem here.

%==============================================================================
\section{Empirical Validation}\label{sec:empirical}
%==============================================================================
The following offline checks and small live panels examine the assumptions
and operational context of the three theorems and the controller pilot
under a pre-registered protocol (decision rules fixed before scoring; see
\path{experiments-next/README.md}). The measurements test assumptions and
operational relevance; they are not used to prove the mathematical statements.

\paragraph{Setup.} One sampling pool: Qwen3-8B, GSM8K train split,
$512$ prompts $\times$ $32$ rollouts at $T{=}1.0$ (Tinker sampling client;
ChatML prompt and binary boxed-answer reward identical to the audit
runners). All analyses are offline resampling of this pool. An initial
billing-truncated pool ($n{=}350$, explicitly stamped) produced identical
verdicts; numbers below are from the complete $n{=}512$ pool. Artifacts:
\texttt{experiments-next/results/*qwen3-8b\_train\_n512\_r32\_s42*.json}.

\paragraph{E-T1: coverage of the ZVF confidence interval.}
\begin{center}
\begin{tabular}{@{}lcc@{}}
\toprule
Batching regime & Estimand & 95\% CI coverage \\
\midrule
shuffled (iid), Wald, $m\ge 64$ & global & 0.93--0.96 \\
shuffled (iid), Wilson, all tested $(m,K)$ & global & 0.95--0.98 \\
difficulty-sorted (curriculum) & global & 0.07--0.22 \\
difficulty-sorted (curriculum) & \emph{local (window)} & 0.944 \\
stratified batch composition & global & 0.996 \\
\bottomrule
\end{tabular}
\end{center}
The Wald interval \eqref{eq:ci} is adequate at $m\ge 64$ and the Wilson
interval is well calibrated everywhere; we adopt Wilson operationally.
The curriculum-ordering ``failure'' is an \emph{estimand mismatch}: the
interval remains calibrated for the local (curriculum-stage) ZVF while
losing the global one, and stratified batch composition restores global
validity (conservatively). This resolves the across-group-i.i.d.\ threat
(Assumption~\ref{ass:across-iid}) into a labeling requirement rather than
an invalidation.

\paragraph{E-T2: the waiting-time quantile.} With $G{=}8$ and
$\delta{=}0.1$, the geometric $(1{-}\delta)$-quantile
\[
N_{\mathrm{disc}}(\theta_G)
=G\left\lceil\frac{\ln(1/\delta)}{-\ln\theta_G}\right\rceil
\]
matched the observed
$(1{-}\delta)$-quantile of rollouts-to-first-informative-group within
0--5\% in all six difficulty strata. In the hardest stratum
($\hat p\approx 0.01$, $\mathrm{ZVF}=0.886$) the quantile is
\textbf{160 rollouts} (ratio 1.00). The geometric waiting-time model is
essentially exact on this reward distribution. Note what is validated:
a \emph{reliability budget} --- the rollouts needed to make a fully
uninformative batch improbable at level $\delta$ --- not a minimum;
informative groups typically arrive much sooner.

\paragraph{E-T3a: the optimal group size (reinterpreted).} The measured
$J$-argmax was 2--3 in every stratum and pooled, agreeing with the
analytic argmax. \textbf{This agreement is now understood to be
algebraically guaranteed, not an empirical test}: by \eqref{eq:J23},
$J(2)=J(3)$ exactly and $G^\star\in\{2,3\}$ for \emph{every} prior, so
the measurement can only confirm the objective's arithmetic (which it
does --- a model-fit sanity check); it cannot validate a data-dependent
prediction, because the objective makes none. What survives is
objective-relative: \emph{if} signal-per-rollout \eqref{eq:efficiency}
is the accepted accounting, prevailing $G\in[8,16]$ practice is
several-fold inefficient per rollout. A first training-side
corroboration: a $G{=}2$ adapter matched a $G{=}16$ adapter's held-out
pass@32 improvement at $8\times$ fewer rollouts per step (single seed;
extended to a two-seed matched-budget panel in E-R2b below, which also
shows the endgame cost of small $G$). Experiment E-T3b under a richer
objective remains the decisive test.

\paragraph{E-B: controller pilot (offline).} Treating stratum selection as
a Thompson-sampling bandit with reward = ``group is informative'' recovers
97\% of the static-oracle's informative groups per 1{,}000 rollouts
(120.7 vs.\ 124.3), beats uniform allocation by 28\%, and exceeds a
hardest-first policy by $8.8\times$ --- difficulty and signal are not the
same axis, exactly as $J(G,p)$ predicts. Caveat: frozen-pool simulation;
the training-time pilot must confirm under the distribution shift induced
by learning itself.

\paragraph{E-R2b: matched-budget training at $G{=}2$ vs.\ $G{=}16$.} Four
LoRA (rank~4) training runs on Qwen3-8B/GSM8K at a matched budget of
2{,}560 rollouts per arm --- $G{=}2\times160$ steps vs.\ $G{=}16\times20$
steps, batch 8, 512-token completions, seeds $\{123,456\}$ --- extend the
single-seed corroboration above. The $G{=}2$ arms convert the budget into
$8\times$ more optimizer steps and drive train reward to
$\approx0.9$--$1.0$ on the sampled pool --- small-$G$ per-rollout
efficiency in action --- before finishing in sustained
$\ZVF\approx0.75$--$1.0$: the all-\emph{correct} zero-variance wall, the
mirror image of T2's all-incorrect wall. The $G{=}16$ arms end
mid-learning (train reward $\approx0.3$--$0.5$ after their 20 steps) with
$\ZVF$ at $0$--$0.25$ throughout and signal intact. Small $G$ buys more optimizer steps per rollout early, then
exhausts its own gradient signal as $\hat p\to 1$; group-size policy is
therefore a \emph{schedule} question, exactly the regime the controller of
Section~\ref{sec:controller} targets. Artifacts:
\texttt{tinker-runs/results/er2b\_*.json}.
This comparison confounds group size with an eightfold difference in optimizer
steps and reports training-pool telemetry, not a held-out causal group-size
effect. Two seeds cannot establish generalization or a preferred schedule.

\paragraph{E-P4: advantage normalization and the length-bias premise.}
A completion-capped (1{,}024-token) six-arm panel --- GRPO vs.\ Dr.GRPO's
$/\sigma$-free advantage \cite{tong2025drgrpo}, three seeds each, 30
steps, $G{=}8$, batch 4, Qwen3-8B/GSM8K --- tests the premise that
std-normalized advantages inflate completion length. Neither loss
inflates length: mean completion length \emph{declines} 3.8--12.2\% over
training in all six arms (GRPO $1004{\to}905$, $981{\to}944$,
$996{\to}900$; Dr.GRPO $999{\to}931$, $972{\to}902$, $1000{\to}878$
tokens). Nor does the $/\sigma$ removal show a ZVF signature: late-run
ZVF is 0.45/0.70/0.72 (Dr.GRPO) vs.\ 0.47/0.47/0.55 (GRPO). At this
scale the length-bias motivation for the normalization change has no
observable footprint on either axis (Tier~C, $n{=}3$ seeds per loss).
The 1,024-token ceiling is close to the reported early means, so censoring can
hide length growth. Batch-level ZVF is coarsely quantized and three seeds have
low power; failure to see a signature does not refute a length-bias mechanism.
Provenance: an earlier version of this panel was invalidated by an
unwired loss flag --- both ``arms'' had trained vanilla GRPO --- and was
caught only by reading the runner, not from any output-level signal; the
invalid artifacts are preserved under explicit
\texttt{.invalid\_actually\_grpo.json} names and the panel was rerun
(2026-07-11). The companion survival-audit paper turns this class of
failure into protocol.

\paragraph{Threats to validity.} Single base model and task family
(Qwen3-8B, GSM8K); binary rewards only; $\hat p$ plug-in bias at $R{=}32$
rollouts per prompt; and training-side corroborations use two seeds per
arm (E-R2b) or three per loss (E-P4).
These measurements do not establish the assumptions on other models, tasks, or
reward types.

%==============================================================================
\section{Discussion and Limitations of the Theory}\label{sec:discussion}
%==============================================================================
\paragraph{What is conditionally established.} Under i.i.d.\ groups and binary
verifiable rewards: (T1) $\ZVF_t$ is an unbiased, asymptotically normal
estimator of a well-defined population quantity with a closed-form CI; (T2)
high population \ZVF{} implies a diverging \emph{reliability budget},
with $N_{\mathrm{disc}}(\theta_K)\sim K\ln(1/\delta)/(1-\theta_K)$ as
$\theta_K\to1^-$ --- the $(1{-}\delta)$-quantile of rollouts to a potentially nonzero centered
reward-contrast contribution under the stated assumptions; (T3) the
signal-per-rollout objective has a well-defined optimum, which is the
\emph{universal} $G^\star\in\{2,3\}$ (eq.~\ref{eq:J23}), prior-independent.

\paragraph{What is not established.}
\begin{itemize}[leftmargin=1.5em,itemsep=1pt]
  \item \textbf{T1 is a binomial-proportion result, not a new U-statistic
        theorem.} A different degree-2 estimand would require a separate
        degeneracy analysis; no such claim is made here.
  \item \textbf{T2 bounds rollouts to a potentially nonzero centered
        reward-contrast contribution under the stated assumptions, not
        rollouts to improvement or a total gradient.} Corollary~\ref{cor:T2-traj} supplies an
        exact-coverage lower bound when the population $\ZVF$ is estimated.
  \item \textbf{T3's objective $S(p,G)$ is a modeling choice,} and under it
        the optimum is the universal, prior-independent $\{2,3\}$
        (eq.~\ref{eq:J23}) --- so T3 as stated supports no data-adaptive
        group-size policy; a richer objective and prospective experiment are
        required for any adaptive-$G^\star$ claim.
  \item \textbf{Across-group i.i.d.\ (Assumption~\ref{ass:across-iid})
        underlies all three} and is violated by curriculum, replay, and
        within-epoch correlation. \emph{Resolved empirically (E-T1):} under
        curriculum ordering the CI remains calibrated for the \emph{local}
        (curriculum-stage) ZVF --- the quantity a stage-local controller
        acts on --- while global inference requires stratified batch
        composition. The threat is thus an estimand-labeling requirement,
        not an invalidation.
\end{itemize}

\paragraph{Empirical status of the theorems.} Superseded by
Section~\ref{sec:empirical}, which reports the complete $n{=}512$
validation (T1 coverage and estimand resolution, T2 floor tightness
1.00--1.05, T3 $G^\star\in\{2,3\}$ agreement, controller pilot). The
truncated-pool ($n{=}350$) artifacts are retained alongside as the
pre-registration record; no verdict changed between $n{=}350$ and
$n{=}512$.

\paragraph{No empirical claims beyond the cited artifacts.} This draft
reports six substudies in Section~\ref{sec:empirical}, including offline
resampling and small live training panels; it does not report a new campaign
beyond those checked-in artifacts. Numeric observations are restricted to the
named artifacts and are not independent validations of algebraic identities.

%==============================================================================

\bibliographystyle{plainnat}
\bibliography{../../platform_hybrid/paper/references}
\end{document}
